% Statistical inference chapter skeleton. Requires \TemplatePrompt.
\section{统计推断模型与证据}

\subsection{估计目标与识别条件}
\TemplatePrompt{定义总体、样本、估计量或效应量、零假设、备择假设及识别所需条件。}

\subsection{描述性基线}
\TemplatePrompt{使用原始组均值、比例、截距模型或未调整效应作为基线，并与调整后估计保持同一量纲。}

\subsection{模型、协变量与抽样结构}
\TemplatePrompt{给出似然或估计方程、协变量选择、分层/聚类结构、缺失机制和参数单位。}

\subsection{假设诊断与估计方法}
\TemplatePrompt{检查分布、残差、方差、独立性、多重比较和校准；说明稳健标准误或重采样设置。}

\subsection{效应量、不确定性与基线比较}
\TemplatePrompt{同时报告效应量、置信区间、样本量和基线差异，不以显著性替代实际意义。}

\subsection{稳健性与结论边界}
\TemplatePrompt{检验替代模型、异常点、子样本和未观测混杂风险；明确相关、预测与因果结论的边界。}
